\chapter*{Conclusion}
\addcontentsline{toc}{chapter}{Conclusion}
\vspace{-0.5cm}
%L'objectif de ce stage était d'implémenter une méthode d'\textbf{interpolation linéaire} dans \textbf{Antares} puis de comparer sa précision avec la \textbf{méthode IDW} dans le cas de propagation acoustique avec l'analogie FWH, tout en optimisant la vitesse d'exécution du code.
%Premièrement, une recherche documentaire sur les méthodes d'interpolation implémentables dans notre cas a été menée.
%Ensuite, la méthode linéaire a été implémentée en utilisant une méthode particulière pour les cellules non triangulaires ou rectangulaires.
%L'efficacité a été augmentée dans les cas multi-instants partagés en évitant le recalcule des coefficients d'interpolation à chaque itération. Ces mêmes coefficients peuvent être récupérés par l'utilisateur qui voudrait appeler deux fois le traitement pour une base ayant une même structure.
%La méthode linéaire a été testée sur des cas simples et industriels. Elle fonctionne correctement sur tous les types de cellules, sauf dans certains cas très complexes comme des prismes ayant des faces opposées de tailles différentes et/ou non parallèles.
%Dans le cas de l'aéroacoustique, la méthode linéaire est plus précise que IDW dans tous les cas testés. Cependant, des résultats expérimentaux ont montré que des coefficients (\(N\), \(p\)) autour de (10, 10) donnaient de très bons résultats. Dans ce cas, IDW peut s'avérer plus précis que la méthode linéaire.


%The goal of this internship was to optimise the numerical performance of a centrifugal pump using DA.
%After preparing the experimental data and adapted the OpenFoam solver and CONES code, data was able to be assimilate by the EnKF.
%First results was not convincing, because a recirculation bubble still present in the diffuser. Few more implementation didn't allow the reciculation bubble to break.
%Then, ...


\textbf{Technical and human review}
% ///// Reprendre les objectifs définis à la fin de la section 2.1 //////

%Bilan technique et humain. Liez chaque aspect technique à un aspect humain. Ne traitez pas tout ce qui est technique d’un côté et tout ce qui relève de l’humain d’un autre. Réflexion personnelle sur les spécificités de l’entreprise, du secteur d’activité et du métier que vous avez découverts.

The first thing that I discovered during this internship is the laboratory environment, which is different from a company. The work is more focused on research, and there is a lot of freedom to explore new ideas and methods. The team is very collaborative and supportive, and I learned a lot from my colleagues. I also discovered the importance of communication and documentation in a research environment, as it is essential to share knowledge and results with others. There are still similaritys with a company, such as the social and environmental responsibility, that I presented in a first part.

After having done a state of art on the topics, we saw that the DA method is a good candidate to DNS or costly LES, for industrial applications, because it allows to have a good accuracy with a lower computational cost. I saw that an international expertise on application relies on this topic has been developed in the last few years on the laboratory. Numerical ingredients and usefull libraries were presented. The state of art was a good opportunity to learn more about the DA method and its applications. It was not too hard because most of information were already known by the laboratory. I also learned a lot about the IBM method and its applications, which is a good alternative to the traditional methods for the simulation of complex geometries.

Then, an explanation of the pump modeling, the pre-processing of the experimental data and the implementation of the IBM method in OpenFOAM were given. The redaction of this part helps me to structure my thoughts and to better understand the different steps of the pre-processing. I also learned a lot about the IBM method and its implementation in OpenFOAM, which was not easy at the beginning because I had to learn a lot about OpenFOAM and C++ programming. The implementation of the IBM method was a good opportunity to learn more about OpenFOAM and its functionalities.

Finally, we presented the first results of the numerical simulation and the first application of the DA method on the centrifugal pump case.
The results show that the IBM method as applied here does not give us a result consistent with physical observations and the DA does not allow us to improve this result. The solution is not converging and the velocity field is not physical, with a recirculation bubble in the diffuser. Several strategies were tested to try to improve the solution, but none of them allowed to obtain a satisfactory result.

I really enjoyed this internship beacause I learned a lot about all the topics that I love. That is why I decided to continue to work on this topic for my PhD thesis "Combining experimental and numerical techniques for the analysis of axial compressor internal flows using Data Assimilation" in the same laboratory.

%\subsection*{Technical and human review}

\textbf{Outlook}

To continue this work, it would be interesting to try to change the IBM method (try the "interpolation method" for example) or to try an instationnary (URANS or LES) simulation to see if it can overcome these non-physical results.

The DA method is still not widely used in the industry, but it has a lot of potential for the future. The main challenge is to adapt the DA method to the specific needs of the industry, such as the use of complex geometries and the need for real-time applications. I discovered that DA can be long to implement and sometimes difficult to converge. At same time, EnKF needs a huge number of CPUs. Sometimes, I need to wait several days to launch a simulation. It would be interesting to try to reduce the computational cost of the DA method, by using more efficient algorithms or by using machine learning techniques, as presented in the state of art.
The IBM method is a good alternative to the traditional methods for the simulation of complex geometries, because it allows to have a good accuracy with a lower computational cost. An other challenge is to improve the accuracy of the IBM method.

\begin{comment}
J’ai été enthousiasmé par ce stage au CERFACS, à plusieurs titres. Il m’a non seulement permis de développer et renforcer de nombreuses connaissances dans le domaine informatique mais également sur le plan humain. Il m’a donné une belle occasion de découvrir la vie d’un laboratoire de recherche.

L'objectif de ce stage était d'implémenter une méthode d'interpolation linéaire dans Antares puis de comparer sa précision avec la méthode IDW dans le cas de propagation acoustique avec l'analogie FWH, tout en optimisant la vitesse d'exécution du code.
Après avoir exploré la structure d'Antares, une recherche documentaire a été menée pour savoir quelle méthode, implémentable dans la librairie, pouvait être candidate pour compléter la méthode IDW et linéaire. Quelques méthodes polynomiales et géostatistiques ont pu être mises en évidence. Une présentation un plus détaillée de la méthode par voisin le plus proche, IDW et linéaire a aussi été réalisée.
L'un des principaux résultats de ce stage a été l'implémentation globalement réussie de la méthode d'interpolation linéaire, qui s'est révélée  généralement plus précise que la méthode de Pondération Inverse à la Distance (IDW) précédemment codée.
Coder l’interpolation linéaire jusqu’en 3D n’a été ni très difficile, ni chronophage. Cependant, l’intégrer au code déjà existant et faire les tests fut plus délicat. Le soutien de mon maître de stage a été primordial pour me montrer comment fonctionnent les outils au CERFACS et m'aider à sortir d'impasses.
Cette nouvelle méthode a été testée sur des cas académiques mais aussi sur des cas industriels tels que sur une vue en coupe à la position (0, 0, 0.01), selon l'axe z de la chambre de combustion preccinsta.
L'étude sur les coefficients \(N\) et \(p\) de la méthode IDW a montré que pour des paramètres (\(N\), \(p\)) autour de (10, 10) nous pouvons obtenir de meilleurs résultats qu'avec la méthode linéaire, dans le cas d'aéroacoustique étudié.
Ce travail a aussi amélioré les performances du code (cinq fois plus rapide pour la méthode IDW dans un cas test d'aéroacoustique).
Le code ajouté pourrait aussi faciliter l'intégration de méthodes d'ordre supérieur grâce à certaines fonctions python réutilisables.
J'ai ressenti tout au long de mon stage l'importance de réaliser un travail ouvrant des perspectives nouvelles pour un chercheur.
Enfin, ce stage n'a pas été qu'une expérience professionnelle enrichissante. Il a conforté mon intérêt pour les méthodes numériques appliquées et la recherche en calcul scientifique. Je suis particulièrement reconnaissant pour l'accompagnement dont j'ai bénéficié tout au long de cette expérience et pour les opportunités de développement personnel et professionnel qu'il m'a offert.
\end{comment}

%\subsection*{Outlook}
%Pour continuer à améliorer le traitement d'interpolation d'Antares, on pourrait implémenter l'une des méthodes d'ordre supérieur citées dans la section \ref{s223}.
%On pourrait également compléter les tests de la méthode linéaire sur la chaîne aéroacoustique, en observant la polaire de l'amplitude en fonction de l'observateur dans le cas d'un Mac non nul.
%Afin d'améliorer la rapidité, le code pourrait être passé en parallèle pour pouvoir s’exécuter sur plusieurs processeurs à la fois.
%Une petite amélioration qui permettrait de rendre le code compatible avec plus de bases serait d'inclure les bases "multizones" à l'interpolation linéaire.
% Faire la ref


%\subsection*{Personal reflection}


%Le CERFACS est un laboratoire privé à la pointe des simulations numériques. Il s'est d'abord spécialisé dans l'aérodynamique et s’intéresse maintenant de plus en plus à la combustion. Il adapte ses domaines de recherche aux domaines importants de l'industrie.
%Durant ce stage, au-delà des aspects techniques, j'ai eu un bel aperçu du métier de chercheur tant au travers de mes missions que dans mes rencontres. Il s'agit de rechercher et de s'appuyer sur le travail déjà réalisé autour de notre problématique.
%Ensuite, il faut explorer des pistes qui peuvent s'avérer infructueuses après les avoir longuement explorées.
%Enfin, un travail de rédaction est nécessaire pour expliquer clairement notre travail aux scientifiques.